\chapter{Undecidable Problems in Mathematics}\label{ch:undecidable}

\begin{center}\itshape
Rewriting words, matching dominoes, tiling floors,\\
iterating fractions, solving equations in whole numbers:\\
each of them can run a Turing machine, so none of them has an algorithm.\\
\upshape(the content of this chapter, in our words)
\end{center}

On 8 August 1900, at the International Congress of Mathematicians in Paris, David Hilbert spoke about the problems that the new century should solve. He read out only ten of them; the printed version had twenty-three. The tenth, which he did not have time to read, takes three lines: ``Given a diophantine equation with any number of unknown quantities and with rational integral numerical coefficients: to devise a process according to which it can be determined by a finite number of operations whether the equation is solvable in rational integers.'' Hilbert did not imagine that the answer might be that no such process exists; there was not yet any precise notion of a process to prove it.

Thirty-six years later there was one (Chapter~\ref{ch:turing}), and the suspicion grew. Around 1950 Martin Davis, who had learnt logic from Emil Post, conjectured that the sets of numbers definable by Diophantine equations are exactly the recognisable ones, which would make the tenth problem unsolvable. Julia Robinson, at Berkeley, was attacking the same problem from the other side, trying to express exponentiation by a polynomial equation. In 1961 Davis, Hilary Putnam and Robinson proved the conjecture for equations in which exponents may also be unknowns. One lemma was missing: a single polynomial equation whose solutions grow exponentially. Robinson had shown how to use one, but nobody could find one.

Robinson later told Constance Reid that each year, when she blew out the candles on her birthday cake, she wished that the tenth problem would be solved, not that she would solve it herself: ``I felt that I couldn't bear to die without knowing the answer.'' On 15 February 1970 Davis telephoned her from New York: John Cocke had just returned from Moscow with the news that Yuri Matiyasevich, a twenty-two-year-old mathematician in Leningrad, had found the equation, using the Fibonacci numbers. Robinson wrote to him: ``If you really are 22, I am especially pleased to think that when I first made the conjecture you were a baby and I just had to wait for you to grow up!'' The four names are now joined in the \emph{MRDP theorem}: there is no algorithm for the tenth problem.

This chapter is a gallery of such problems, from algebra, combinatorics, geometry and number theory. They look nothing like Turing machines, but each of them can simulate one, and that is enough to make it undecidable.

\begin{think}
You are given a set of square tiles with coloured edges, as many copies of each as you like, and must cover an infinite floor so that touching edges have the same colour. If you succeed on a $10\times10$ square, then a $100\times100$ square, how would you ever know whether the whole floor can be covered? And if the tiles could only cover the floor in a pattern that never repeats, could you still find out?
\end{think}

\section{Rewriting words}\label{sec:thue}

In 1914 the Norwegian mathematician Axel Thue studied a game with words: you may replace certain factors by others, according to a fixed list of rules, and you ask whether one word can be turned into another.

\begin{definition}\label{def:semithue}
A \emph{semi-Thue system} $R$ over a finite alphabet $\Sigma$ is a finite set of rules $u\to v$ with $u,v\in\Sigma^+$ (nonempty words). We write $xuy\Rightarrow xvy$ for every rule $u\to v$ and all words $x,y$, and $\Rightarrow^*$ for the reflexive and transitive closure. A word $z$ is a \emph{successor} of $w$ if $w\Rightarrow z$. The \emph{Thue system} of $R$ allows the rules in both directions; we write $\Leftrightarrow^*$ for its derivations, an equivalence relation compatible with concatenation. The quotient $\Sigma^+/\Leftrightarrow^*$ is the semigroup \emph{presented} by the generators $\Sigma$ and the relations $u=v$; two words represent the same element if and only if $x\Leftrightarrow^*y$. The \emph{word problem} asks, given $x$ and $y$, whether they are equal there.
\end{definition}

Thue asked for an algorithm deciding $x\Rightarrow^*y$. A semi-Thue system is a very simple programming language, and a Turing machine fits into it.

\begin{theorem}[Markov, 1947; Post, 1947]\label{thm:semithue}
For every deterministic Turing machine $M$ there is a semi-Thue system $R_M$, computable from $M$, with a letter $H$ such that for every input $w$: $[q_0w]\Rightarrow^*H$ if and only if $M$ accepts $w$. Moreover $R_M$ has a set of \emph{marked} letters such that both sides of every rule contain exactly one marked letter, and every word with exactly one marked letter has at most one successor.
\end{theorem}

\begin{proof}
Let $M$ have states $Q$, tape alphabet $\Gamma$ with blank $\sqcup$, and transition function $\delta$ (Definition~\ref{def:tm}). The alphabet of $R_M$ is $\Gamma\cup Q\cup\{[,],p,H\}$, the new letters being distinct from the others, and the marked letters are those of $Q\cup\{p,H\}$. The word $[uqv]$ with $u,v\in\Gamma^*$ stands for the configuration in which the tape holds $uv$ (blanks elsewhere), the state is $q$ and the head is on the first letter of $v$, or on the blank to the right of $u$ if $v$ is empty. For each state $q\notin\{q_{\mathrm{acc}},q_{\mathrm{rej}}\}$ and letter $a\in\Gamma$, with $\delta(q,a)=(r,b,D)$, the rules are:
\begin{itemize}
\item if $D=\mathrm R$: $qa\to br$, and, when $a=\sqcup$, also $q]\to br]$;
\item if $D=\mathrm L$: $cqa\to rcb$ for every $c\in\Gamma$ and $[qa\to[r\sqcup b$; when $a=\sqcup$, also $cq]\to rcb]$ for every $c\in\Gamma$ and $[q]\to[r\sqcup b]$.
\end{itemize}
They perform one step of $M$, the rules with $]$ for a head on the blank right of the written part, the rules with $[$ for a head that moves left of it. The \emph{clean-up} rules are $q_{\mathrm{acc}}a\to q_{\mathrm{acc}}$ for $a\in\Gamma$, $q_{\mathrm{acc}}]\to p]$, $ap\to p$ for $a\in\Gamma$, and $[p]\to H$; they erase everything once $M$ accepts. There are no rules for $q_{\mathrm{rej}}$.

Each side of each rule contains exactly one marked letter. In a word with exactly one marked letter, a rule can only be applied around it, and its neighbours decide which: for a state $q\notin\{q_{\mathrm{acc}},q_{\mathrm{rej}}\}$, the letter after $q$ (a letter $a$ or $]$) fixes the transition, and if it moves left the letter before $q$ (a letter $c$ or $[$) fixes the rule; for $q_{\mathrm{acc}}$ the next letter, for $p$ the previous one, decides; $H$ and $q_{\mathrm{rej}}$ have no rules. So such a word has at most one successor.

Starting from $[q_0w]$, the successors follow the computation of $M$ on $w$, one rule per step, through words $[uqv]$ describing the successive configurations. If $M$ accepts, we reach $[uq_{\mathrm{acc}}v]$, then $[uq_{\mathrm{acc}}]$, $[up]$, $[p]$ and $H$. If $M$ rejects, we reach a word with $q_{\mathrm{rej}}$, which has no successor; if it runs forever, so does the derivation. Since the derivation from $[q_0w]$ is unique, and $H$ can only be produced from $[p]$, which can only come from $q_{\mathrm{acc}}$, we get $[q_0w]\Rightarrow^*H$ if and only if $M$ accepts $w$.
\end{proof}

In the semi-Thue system the rules go one way, like time. In a semigroup the relations go both ways, and a derivation may wander backwards. Post noticed that for deterministic systems this does not help.

\begin{lemma}[Post, 1947]\label{lem:post}
Let $W$ be a set of words closed under applying rules of $R$ in both directions, in which every word has at most one successor, and let $H\in W$ have no successor. Then for $x\in W$: $x\Leftrightarrow^*H$ if and only if $x\Rightarrow^*H$.
\end{lemma}

\begin{proof}
Take a chain $x=w_0,w_1,\dots,w_n=H$ of minimal length in which each step is a rule applied forwards ($w_i\Rightarrow w_{i+1}$) or backwards ($w_{i+1}\Rightarrow w_i$); all $w_i$ lie in $W$. If some step is backwards, consider the last one, $w_{i+1}\Rightarrow w_i$. It is not the last step, since $H$ has no successor, so the next step is forwards: $w_{i+1}\Rightarrow w_{i+2}$. As $w_{i+1}$ has only one successor, $w_i=w_{i+2}$, and deleting $w_{i+1},w_{i+2}$ shortens the chain, a contradiction. So all steps are forwards.
\end{proof}

\begin{corollary}[Markov, 1947; Post, 1947]\label{cor:semigroup}
There is a finitely presented semigroup whose word problem is undecidable.
\end{corollary}

\begin{proof}
Apply Theorem~\ref{thm:semithue} to the universal machine $U$ of Theorem~\ref{thm:universal}, which accepts $\langle M,w\rangle$ if and only if $M$ accepts $w$. The words with exactly one marked letter form a set $W$ closed under the rules in both directions, and the lemma applies: $[q_0\langle M,w\rangle]=H$ in the semigroup presented by $R_U$ if and only if $[q_0\langle M,w\rangle]\Rightarrow^*H$, if and only if $M$ accepts $w$. An algorithm for the word problem would decide $A_{\mathrm{TM}}$, contradicting Theorem~\ref{thm:halting}.
\end{proof}

Grigori Tseitin found in 1958 a semigroup with five generators $a,b,c,d,e$ and seven relations,
\[
ac=ca,\quad ad=da,\quad bc=cb,\quad bd=db,\quad eca=ce,\quad edb=de,\quad cca=ccae,
\]
whose word problem is undecidable, a result we quote. For groups, where every generator has an inverse, the same theorem was proved by Pyotr Novikov (1955) and William Boone (1958) with much more work, and it is the starting point of the undecidability results of topology: no algorithm decides whether two four-dimensional manifolds, given by triangulations, are homeomorphic (Markov, 1958).

\section{Post's correspondence problem}\label{sec:pcp}

Post found in 1946 a problem with no machines and no rewriting in sight, which has become the standard tool for proving undecidability.

\begin{definition}\label{def:pcp}
An instance of \emph{Post's correspondence problem} (PCP) is a finite list of \emph{dominoes} $(t_1,b_1),\dots,(t_k,b_k)$, pairs of nonempty words over an alphabet $\Sigma$. A \emph{match} is a nonempty sequence of indices $i_1,\dots,i_m$ with $t_{i_1}t_{i_2}\cdots t_{i_m}=b_{i_1}b_{i_2}\cdots b_{i_m}$: the words on top and at the bottom of the dominoes, laid side by side, are the same. In the \emph{modified} problem (MPCP), a match must start with the first domino.
\end{definition}

Figure~\ref{fig:pcp} shows an instance and a match. A match can be searched for by trying all sequences, which proves that the set of instances with a match is recognisable; the difficulty is to know when to stop.

\begin{figure}[ht]
\centering
\fig{ch04-pcp}
\caption{An instance of Post's correspondence problem with four dominoes (top), and a match $2,1,3,2,4$ (bottom): both rows spell $abcaaabc$, cut in different places. Exercise~\ref{ex:pcphand} asks for a match by hand in another instance; the notebook searches for matches by computer.}
\label{fig:pcp}
\end{figure}

\begin{theorem}[Post, 1946]\label{thm:pcp}
Post's correspondence problem is undecidable.
\end{theorem}

\begin{proof}
\emph{Step 1: MPCP.} Let $R=R_U$ be the system of Theorem~\ref{thm:semithue} for the universal machine, and $x=[q_0\langle M,w\rangle]$. Every word derivable from $x$ is either $H$ or of the form $[z]$ with $z$ free of brackets and of $H$, since every rule preserves this form except $[p]\to H$, which needs the whole word. So the only word derivable from $x$ that contains $H$ is $H$. Let $\#$ be a new letter, and consider the dominoes
\[
D_1=(\#,\ \#x\#),\qquad (\#,\#),\qquad F=(H\#\#,\ \#),
\]
\[
(a,a)\ \text{for each letter }a\text{ of }R,\qquad (u,v)\ \text{for each rule }u\to v.
\]
We show that this MPCP instance has a match if and only if $x\Rightarrow^*H$, that is, if and only if $M$ accepts $w$.

If $x=x_0\Rightarrow x_1\Rightarrow\cdots\Rightarrow x_n=H$, start with $D_1$: the top reads $\#$ and the bottom $\#x_0\#$. For $j=0,\dots,n-1$, lay the dominoes $(a,a)$ under the letters of $x_j$, except for the one factor $u$ that is rewritten, which gets $(u,v)$, and then $(\#,\#)$: the top grows by $x_j\#$ and the bottom by $x_{j+1}\#$. The top is now $\#x_0\#\cdots\#x_{n-1}\#$ and the bottom $\#x_0\#\cdots\#x_{n-1}\#H\#$; the domino $F$ adds $H\#\#$ to the top and $\#$ to the bottom, and both rows read $\#x_0\#\cdots\#x_{n-1}\#H\#\#$.

Conversely, let $s$ be the common word of a match, and call \emph{segments} of $s$ the factors $\sigma_1,\sigma_2,\dots$ between consecutive letters $\#$, so that $s=\#\sigma_1\#\sigma_2\#\cdots$. The domino $F$ is used, since only $F$ has more $\#$ on top than at the bottom while $D_1$ has fewer. Let $E_1=D_1,E_2,\dots,E_r$ be the dominoes containing $\#$ that come before the first $F$, and call the others \emph{letter dominoes}. On top, $E_i$ contributes the $i$th letter $\#$ of $s$, so the letter dominoes between $E_i$ and $E_{i+1}$ have tops forming $\sigma_i$ (and those between $E_r$ and $F$, followed by the $H$ of $F$, form $\sigma_r$). At the bottom, after $E_i$ there are $c_i\ge i+1$ letters $\#$, because $D_1$ contributes two; so the bottoms of the letter dominoes between $E_i$ and $E_{i+1}$ form the segment $\sigma_{c_i}$, and $\sigma_i\Rightarrow^*\sigma_{c_i}$ by applying their rules one at a time. Every segment written at the bottom before $F$ is therefore $x$ (written by a copy of $D_1$) or is derived from an earlier segment. By induction, all segments $\sigma_k$ with $k<c_r$ are derivable from $x$, and in particular $\sigma_r$, since $r<c_r$. But $\sigma_r$ ends with $H$, so $\sigma_r=H$ and $x\Rightarrow^*H$.

\emph{Step 2: from MPCP to PCP.} For a word $u=a_1\cdots a_n$ write $\star u=\star a_1\star a_2\cdots\star a_n$ and $u\star=a_1\star a_2\star\cdots a_n\star$, with a new letter $\star$. To the instance $(t_1,b_1),\dots,(t_k,b_k)$ of MPCP associate the PCP instance
\[
(\star t_1,\ \star b_1\star),\qquad(\star t_i,\ b_i\star)\ \ (1\le i\le k),\qquad(\star\diamond,\ \diamond),
\]
with a new letter $\diamond$. A match of the first gives a match of the second: replace the first domino by $(\star t_1,\star b_1\star)$, the others by their starred versions, and end with $(\star\diamond,\diamond)$. Conversely, in a match of the second instance, the first domino is $(\star t_1,\star b_1\star)$, the only one whose two words begin with the same letter. Cut the match just after the first occurrence of $(\star\diamond,\diamond)$: since $\diamond$ appears in both rows only through this domino, its first occurrence is at the same place in both rows, and the cut sequence is still a match. It contains $(\star\diamond,\diamond)$ only at the end, and erasing $\star$ and $\diamond$ from it gives a match of the MPCP instance.

An algorithm for PCP would therefore decide whether $M$ accepts $w$, which is impossible.
\end{proof}

The theorem keeps the promise of Chapter~\ref{ch:grammars} about grammars. Fix dominoes $(t_1,b_1),\dots,(t_k,b_k)$ over $\Sigma$, and new letters $\iota_1,\dots,\iota_k$, the \emph{indices}. Let $G_A$ be the grammar $A\to t_iA\iota_i\mid t_i\iota_i$ ($1\le i\le k$), which generates
\[
L_A=\{t_{i_1}t_{i_2}\cdots t_{i_m}\iota_{i_m}\cdots\iota_{i_2}\iota_{i_1}:m\ge1\},
\]
and $G_B$ the same with the $b_i$. A word of $L_A$ has only one parse tree, because its indices, read backwards, list the rules used. And $L_A\cap L_B\ne\emptyset$ if and only if the instance has a match.

\begin{proposition}\label{prop:pcpcomplement}
The complement of $L_A$ in $(\Sigma\cup\{\iota_1,\dots,\iota_k\})^*$ is context-free, and a grammar for it can be computed from the dominoes.
\end{proposition}

\begin{proof}
For $y=\iota_{i_m}\cdots\iota_{i_1}$ write $h(y)=t_{i_1}\cdots t_{i_m}$, so that $L_A=\{h(y)y:y\in I^+\}$, where $I$ is the set of indices. A word is not in $L_A$ if and only if it is not in $\Sigma^*I^+$, a regular condition, or it is $zy$ with $z\in\Sigma^*$, $y\in I^+$ and $z\ne h(y)$. In the second case, either $h(y)$ is a proper prefix of $z$, or there is a first $j$ such that $z=t_{i_1}\cdots t_{i_{j-1}}z'$ where $z'$ does not begin with $t_{i_j}$. The words of the first kind are generated by $N_1\to t_iN_1\iota_i\mid t_iP\iota_i$, where $P$ generates $\Sigma^+$; the words of the second kind by $N_2\to t_iN_2\iota_i\mid Z_iY\iota_i$ ($1\le i\le k$), where $Z_i$ generates the regular language of words of $\Sigma^*$ that do not begin with $t_i$ and $Y$ generates $I^*$. Regular languages are context-free (Chapter~\ref{ch:grammars}), and so is a union of context-free languages.
\end{proof}

\begin{corollary}\label{cor:cfgundecidable}
No algorithm decides, for context-free grammars $G,G'$ over an alphabet $\Sigma$:
\textup{(a)} whether $G$ is ambiguous; \textup{(b)} whether $L(G)\cap L(G')=\emptyset$; \textup{(c)} whether $L(G)=\Sigma^*$; \textup{(d)} whether $L(G)=L(G')$.
\end{corollary}

\begin{proof}
For a PCP instance, form the grammars above. (a) The grammar $S\to A\mid B$, with the rules of $G_A$ and $G_B$, is ambiguous if and only if some word has a derivation through $A$ and one through $B$, that is, if and only if $L_A\cap L_B\ne\emptyset$, if and only if the instance has a match. (b) $L(G_A)\cap L(G_B)=\emptyset$ if and only if there is no match. (c) By Proposition~\ref{prop:pcpcomplement} there is a grammar $G$ for $\overline{L_A}\cup\overline{L_B}$, which is everything if and only if $L_A\cap L_B=\emptyset$. (d) Compare $G$ with a grammar generating all words. In each case an algorithm would decide PCP.
\end{proof}

\section{Tiling the plane}\label{sec:wang}

In 1961 the logician Hao Wang, working on the decision problem of logic at Bell Labs, reduced a class of formulas to a puzzle about tiles.

\begin{definition}\label{def:wang}
A \emph{Wang tile} is a unit square whose four edges carry colours. A \emph{tile set} is a finite set $T$ of tiles; a \emph{tiling} of a region $X\subseteq\mathbb Z^2$ is a map $\tau:X\to T$ such that adjacent tiles agree: the right colour of $\tau(x,y)$ is the left colour of $\tau(x+1,y)$ and the top colour of $\tau(x,y)$ is the bottom colour of $\tau(x,y+1)$, whenever both cells are in $X$. Tiles may be reused but not rotated. The \emph{domino problem} asks whether $T$ tiles the plane $\mathbb Z^2$; the \emph{origin-constrained} problem asks, given also $t_0\in T$, whether some tiling of the plane has $\tau(0,0)=t_0$.
\end{definition}

An infinite object is hard to search for, but compactness reduces it to finite ones.

\begin{proposition}\label{prop:compactness}
A tile set $T$ tiles the plane with $t_0$ at the origin if and only if, for every $n$, it tiles the square $[-n,n]^2$ with $t_0$ at the origin. Consequently, the set of pairs $(T,t_0)$ for which no such tiling exists is recognisable, and so is the set of tile sets that do not tile the plane.
\end{proposition}

\begin{proof}
A tiling of the plane restricts to the squares. Conversely, let $S_n$ be the finite nonempty set of tilings of $[-n,n]^2$ with $t_0$ at the origin; restricting a tiling in $S_{n+1}$ gives one in $S_n$. Call $\tau\in S_n$ \emph{good} if, for every $m\ge n$, it is the restriction of some tiling in $S_m$. The unique element of $S_0$ is good. If $\tau\in S_n$ is good, then for each $m>n$ one of the finitely many extensions of $\tau$ in $S_{n+1}$ is the restriction of a tiling in $S_m$; one of them works for infinitely many $m$, hence for all $m>n$, and is good. So there are good tilings $\tau_0\subset\tau_1\subset\cdots$ of growing squares, and their union tiles the plane. This is König's lemma (Tome~C1, Chapter~2). For the last sentence, examine the squares $[-n,n]^2$ one after the other by exhaustive search, and accept as soon as one cannot be tiled.
\end{proof}

\begin{theorem}\label{thm:wang}
No algorithm decides, given a tile set $T$ and $t_0\in T$, whether $T$ tiles the plane with $t_0$ at the origin.
\end{theorem}

\begin{proof}
No algorithm decides whether a machine halts on the blank tape: if one did, given $\langle M,w\rangle$ we could build the machine that writes $w$ and then runs $M$, and decide the halting problem (Corollary~\ref{cor:halt}). So it suffices to build from each machine $M$ a tile set $T_M$ and a tile $t_0$ such that $T_M$ tiles the plane with $t_0$ at the origin if and only if $M$ does not halt on the blank tape. The idea is that row $y$ of the tiling is the configuration after $y$ steps (Figure~\ref{fig:wang}).

The top and bottom colours are the \emph{cell colours} $a\in\Gamma$ (a cell holding $a$) and $(q,a)$ (a cell holding $a$ under the head, in state $q$), plus a colour $\beta$. The left and right colours are $\emptyset$ (nothing crosses), the \emph{arrows} $q{\to}$ and ${\leftarrow}q$ (the head, in state $q$, crosses the edge to the right or to the left), and two colours $\sigma,\sigma'$. The tiles are, writing a tile as (bottom, top; left, right):
\begin{itemize}
\item the \emph{start} tiles $t_0=(\beta,(q_0,\sqcup);\sigma',\sigma)$, $(\beta,\sqcup;\sigma,\sigma)$, $(\beta,\sqcup;\sigma',\sigma')$, and the \emph{filler} $(\beta,\beta;\emptyset,\emptyset)$;
\item the \emph{copy} tiles $(a,a;\emptyset,\emptyset)$ for $a\in\Gamma$;
\item for each non-halting $q$ and $a\in\Gamma$ with $\delta(q,a)=(r,b,D)$, the \emph{action} tile $((q,a),b;\emptyset,r{\to})$ if $D=\mathrm R$, $((q,a),b;{\leftarrow}r,\emptyset)$ if $D=\mathrm L$;
\item the \emph{arrival} tiles $(c,(r,c);r{\to},\emptyset)$ and $(c,(r,c);\emptyset,{\leftarrow}r)$ for every state $r$ and $c\in\Gamma$.
\end{itemize}

Suppose $\tau$ tiles the plane with $\tau(0,0)=t_0$. The tile right of $t_0$ has left colour $\sigma$, so it is $(\beta,\sqcup;\sigma,\sigma)$, and so on to the right; to the left, likewise, $(\beta,\sqcup;\sigma',\sigma')$. So the top colours of row $0$ spell the initial configuration $C_0$: state $q_0$ on a blank tape. We show by induction that, as long as $M$ has not halted, the top colours of row $y$ spell the configuration $C_y$ after $y$ steps. Let the bottom colours of row $y+1$ spell $C_y$, with the head at $h$ in a non-halting state $q$ reading $a$. The only tile with bottom $(q,a)$ is the action tile, so $\tau(h,y+1)$ is the action tile; say $D=\mathrm R$. Then $\tau(h+1,y+1)$ has left colour $r{\to}$ and bottom colour a letter $c$, so it is the arrival tile with top $(r,c)$. Every other tile of the row has a letter as bottom colour, so it is a copy tile or an arrival tile; but an arrival tile needs an arrow on one side, and an arrow pointing at it can only come from an action tile, which is at $h$. So the other tiles are copy tiles, and the top colours of row $y+1$ spell $C_{y+1}$. The case $D=\mathrm L$ is symmetric. Finally, if $M$ halts after $y$ steps, row $y+1$ would need a tile with bottom $(q,a)$ for a halting state $q$, and there is none. So $M$ does not halt.

Conversely, if $M$ never halts, put the start tiles in row $0$, the fillers in all rows below, and in row $y\ge1$ the action tile at the head of $C_{y-1}$, the arrival tile where the head goes, and copy tiles elsewhere. All edges match, which is a tiling of the plane with $t_0$ at the origin.
\end{proof}

\begin{figure}[ht]
\centering
\fig{ch04-wang}
\caption{The tiling of Theorem~\ref{thm:wang} for a machine with states $A$ (blue) and $B$ (red), and symbols $0$ (white) and $1$ (yellow): $A0\mapsto(B,1,\mathrm R)$, $A1\mapsto(A,1,\mathrm L)$, $B0\mapsto(A,1,\mathrm L)$, $B1\mapsto(A,1,\mathrm R)$. Each tile is drawn as four triangles carrying its edge colours. Row $0$ (green) starts the machine at the origin, and row $y$ holds the configuration after $y$ steps; the head is passed sideways by arrows. This machine never halts, so the pattern continues upwards forever.}
\label{fig:wang}
\end{figure}

By Proposition~\ref{prop:compactness} the set of $(T,t_0)$ without a tiling is recognisable; by the theorem, it is not decidable. The problem without $t_0$ was the one Wang cared about, and he made an observation that seemed to settle it. Say that a tiling is \emph{periodic} if there is $p\ge1$ with $\tau(x+p,y)=\tau(x,y+p)=\tau(x,y)$ for all $x,y$.

\begin{proposition}[Wang, 1961]\label{prop:wangperiodic}
If every tile set that tiles the plane had a periodic tiling, the domino problem would be decidable.
\end{proposition}

\begin{proof}
A periodic tiling with period $p$ is given by a tiling of the square $[0,p-1]^2$ whose right colours of the last column match the left colours of the first, and whose top colours of the last row match the bottom colours of the first; conversely such a square, repeated, tiles the plane. So, given $T$, search in parallel for $p$ such that such a square exists, and for $n$ such that $[-n,n]^2$ cannot be tiled. Under the hypothesis one of the two searches stops (Proposition~\ref{prop:compactness}), and it tells whether $T$ tiles the plane.
\end{proof}

\begin{theorem}[Berger, 1966]\label{thm:berger}
The domino problem is undecidable. Consequently there are \emph{aperiodic} tile sets, which tile the plane but have no periodic tiling.
\end{theorem}

We quote Berger's theorem, the second sentence being Proposition~\ref{prop:wangperiodic}. Without the tile at the origin, nothing forces the computation to start, and a tiling may contain no computation at all, or infinitely many unfinished ones; Robert Berger, Wang's student, overcame this with a set of $20\,426$ tiles that force a hierarchy of squares of sizes $4,8,16,\dots$, each running a computation inside. Raphael Robinson simplified the construction in 1971, and Emmanuel Jeandel and Michaël Rao found in 2015, by computer search, an aperiodic set of $11$ Wang tiles with $4$ colours, and proved that no aperiodic set has fewer tiles or fewer colours. In 2023 David Smith, a retired print technician and amateur, with Joseph Myers, Craig Kaplan and Chaim Goodman-Strauss, found the ``hat'', a single polygon whose copies (rotated and reflected) tile the plane but never periodically.

\section{Counters, fractions and Collatz}\label{sec:fractran}

Turing machines are already simple, but Marvin Minsky showed in 1961 that a machine with no tape at all, only a few counters, is just as powerful.

\begin{definition}\label{def:counter}
A \emph{counter machine} with $k$ counters $c_1,\dots,c_k\in\mathbb N$ is a finite numbered list of instructions, each of the form $\mathrm{INC}(i,j)$: add $1$ to $c_i$ and go to instruction $j$; $\mathrm{DEC}(i,j,l)$: if $c_i>0$, subtract $1$ and go to $j$, otherwise go to $l$; or $\mathrm{HALT}$. It starts at instruction $1$ with given values of the counters.
\end{definition}

\begin{theorem}[Minsky, 1961]\label{thm:minsky}
For every Turing machine $M$ there is a counter machine with three counters, computable from $M$, which, started with $c_1=0$, $c_2=\mathrm{code}(w)$ and $c_3=0$ for a computable coding of the input $w$, halts if and only if $M$ halts on $w$.
\end{theorem}

\begin{proof}
Number the tape symbols $0,\dots,B-1$ with the blank as $0$. A configuration with head at cell $h$ is coded by $L=\sum_{j\ge1}a_{h-j}B^{j-1}$ (the cells to the left of the head, nearest first) and $R=\sum_{j\ge0}a_{h+j}B^j$ (the head cell and those to its right), finite sums since all but finitely many cells are blank; the state is kept in the program. The input $w=a_0a_1\cdots a_{n-1}$ is coded by $\mathrm{code}(w)=\sum a_jB^j$. The machine reads $s=R\bmod B$, writes $s'$, and moves:
\[
\text{right: } (L,R)\mapsto\bigl(BL+s',\ \lfloor R/B\rfloor\bigr),\qquad
\text{left: } (L,R)\mapsto\bigl(\lfloor L/B\rfloor,\ (L\bmod B)+B(s'+B\lfloor R/B\rfloor)\bigr).
\]
With the third counter $T$ as scratch, each piece is a loop. To divide $c$ by $B$: repeat ``decrement $c$ up to $B$ times; if all $B$ decrements succeed, increment $T$'', counting in the instruction number how many succeeded in the last round, which is $c\bmod B$; at the end $c=0$, $T=\lfloor c/B\rfloor$, and we move $T$ back into $c$. To multiply $c$ by $B$: move $c$ into $T$ while adding $B$ to $T$ at each unit. Adding a known digit is a fixed number of increments. Each state and each read symbol gives a block of such instructions; halting states become $\mathrm{HALT}$.
\end{proof}

John Conway turned counter machines into arithmetic. A \emph{FRACTRAN program} is a finite list of positive fractions $f_1,\dots,f_r$. Started on a positive integer $n$, it replaces $n$ by $f_jn$ for the first $j$ such that $f_jn$ is an integer, and repeats; it halts when no $f_jn$ is an integer.

\begin{theorem}[Conway, 1987]\label{thm:fractran}
For every counter machine $C$ there is a FRACTRAN program $F$, computable from $C$, such that $F$ started on $2^{c_1}3^{c_2}\cdots p_k^{c_k}\cdot q_1$ halts if and only if $C$ started with counters $c_1,\dots,c_k$ halts. Consequently, no algorithm decides whether a FRACTRAN program halts on a given integer.
\end{theorem}

\begin{proof}
Assign to counter $i$ the prime $p_i$ (the first $k$ primes) and to instruction $s$ a new prime $q_s$. The number $q_s\prod_ip_i^{c_i}$ codes the machine at instruction $s$ with counters $c_i$. For instruction $s$ put in the list: the fraction $q_jp_i/q_s$ if it is $\mathrm{INC}(i,j)$; the two fractions $q_j/(q_sp_i)$ then $q_l/q_s$, in this order, if it is $\mathrm{DEC}(i,j,l)$; nothing if it is $\mathrm{HALT}$. We may assume that no instruction jumps to itself: otherwise use two copies of the program, the jumps of each copy leading into the other. Then each fraction is in lowest terms, none is $1$, and its denominator contains exactly one $q_s$, so on a code at instruction $s$ only the fractions of instruction $s$ can apply. For $\mathrm{DEC}(i,j,l)$ the first applies exactly when $p_i$ divides the code, that is $c_i>0$, and it decrements $c_i$ and moves to $j$; otherwise the second moves to $l$. So $F$ follows $C$ step by step, and halts exactly at a $\mathrm{HALT}$ instruction. The last sentence follows from Theorem~\ref{thm:minsky} and the undecidability of the halting problem.
\end{proof}

Conway's PRIMEGAME is the program
\[
\tfrac{17}{91},\ \tfrac{78}{85},\ \tfrac{19}{51},\ \tfrac{23}{38},\ \tfrac{29}{33},\ \tfrac{77}{29},\ \tfrac{95}{23},\ \tfrac{77}{19},\ \tfrac1{17},\ \tfrac{11}{13},\ \tfrac{13}{11},\ \tfrac{15}{14},\ \tfrac{15}2,\ \tfrac{55}1.
\]
Started on $2$, it never halts, and the powers of $2$ it meets are $2^2,2^3,2^5,2^7,2^{11},\dots$, the prime powers of two in increasing order, a theorem of Conway that we quote (the notebook watches the first ones appear, after $19$ steps for $2^2$ and $69$ for $2^3$).

A FRACTRAN step multiplies $n$ by a constant depending only on the residue of $n$ modulo $m$, the least common multiple of the denominators, because $f_jn$ is an integer exactly when the denominator of $f_j$ divides $n$. This is the shape of the Collatz map $n\mapsto n/2$ ($n$ even), $3n+1$ ($n$ odd), whose iterates are conjectured to reach $1$ from every start, a problem open since Lothar Collatz posed it in the 1930s.

\begin{corollary}[Conway, 1972]\label{cor:collatz}
No algorithm decides, given $m\ge1$, rationals $a_0,\dots,a_{m-1}$ such that $g(n)=a_{n\bmod m}\,n$ is a positive integer for every positive integer $n$, and a starting value $n_0$, whether the iterates of $g$ starting from $n_0$ reach a fixed point.
\end{corollary}

\begin{proof}
Given a FRACTRAN program from Theorem~\ref{thm:fractran}, none of whose fractions is $1$ (true of the programs of the proof), let $m$ be the least common multiple of the denominators and, for each residue $r$, let $a_r=f_j$ for the first $j$ whose denominator divides $r$ (equivalently, every $n\equiv r$), or $a_r=1$ if there is none. Then $g$ performs the steps of the program, and $g(n)=n$ exactly when the program halts on $n$.
\end{proof}

Stuart Kurtz and Janos Simon proved in 2007 that the question ``does every starting value reach $1$?'', for maps of this kind, is even harder than the halting problem: it is undecidable even for a machine that can solve the halting problem. The Collatz problem itself is one instance; perhaps it is true and unprovable, perhaps not.

\section{Hilbert's tenth problem}\label{sec:hilbert10}

Diophantine equations met us in Tome~1, Chapter~9, where special methods solve special equations. Here we ask what a single polynomial can express.

\begin{definition}\label{def:diophantine}
A set $S\subseteq\mathbb N^k$ is \emph{Diophantine} if there is a polynomial $P(x_1,\dots,x_k,y_1,\dots,y_m)$ with integer coefficients such that
\[
(x_1,\dots,x_k)\in S\iff\exists y_1,\dots,y_m\in\mathbb N:\ P(x_1,\dots,x_k,y_1,\dots,y_m)=0.
\]
\end{definition}

\begin{proposition}\label{prop:dioph}
\textup{(a)} The sets $\{(x,y):x\le y\}$, $\{(x,y):x\mid y\}$, $\{(x,y):x\ne y\}$ and the set of composite numbers are Diophantine. \textup{(b)} Unions, intersections and projections ($\exists$ on a coordinate) of Diophantine sets are Diophantine. \textup{(c)} Every Diophantine set is recognisable.
\end{proposition}

\begin{proof}
(a) Take $y-x-z=0$, $y-xz=0$, $(x-y)^2-z-1=0$ and $x-(y+2)(z+2)=0$, with unknowns $z$ (and $y$ for the last). (b) If $P$ and $Q$ define $S$ and $S'$ with separate unknowns, $PQ=0$ defines $S\cup S'$ and $P^2+Q^2=0$ defines $S\cap S'$; a projection turns a coordinate into an unknown. (c) Given $\bar x$, enumerate all $\bar y\in\mathbb N^m$ (Exercise~\ref{ex:pairing} lists the pairs) and accept when $P(\bar x,\bar y)=0$.
\end{proof}

\begin{theorem}[Matiyasevich, Robinson, Davis, Putnam, 1970]\label{thm:mrdp}
Every recognisable subset of $\mathbb N^k$ is Diophantine.
\end{theorem}

We quote the theorem, and prove below its key number-theoretic lemmas. Two consequences come at once.

\begin{corollary}\label{cor:hilbert10}
There is a polynomial $P(x,y_1,\dots,y_m)$ with integer coefficients such that no algorithm decides, given $x\in\mathbb N$, whether $P(x,\bar y)=0$ has a solution $\bar y\in\mathbb N^m$. In particular Hilbert's tenth problem has no solution, for unknowns in $\mathbb N$ and also in $\mathbb Z$.
\end{corollary}

\begin{proof}
Code the words $\langle M,w\rangle$ by integers; $A_{\mathrm{TM}}$ becomes a recognisable, undecidable set of integers, which is Diophantine by the theorem. For unknowns in $\mathbb Z$, Exercise~\ref{ex:integers} transforms one problem into the other.
\end{proof}

\begin{corollary}[Putnam, 1960; Jones, Sato, Wada, Wiens, 1976]\label{cor:primepoly}
There is a polynomial with integer coefficients whose positive values, for natural values of its variables, are exactly the prime numbers.
\end{corollary}

\begin{proof}
The primes form a decidable set, so by the theorem there is $P$ with $x$ prime $\iff\exists\bar y\ P(x,\bar y)=0$. Let $Q(x,\bar y)=(x+1)\bigl(1-P(x,\bar y)^2\bigr)-1$. If $P(x,\bar y)=0$ then $Q=x$; otherwise $P^2\ge1$ and $Q\le-1$. So the positive values of $Q$ are the primes.
\end{proof}

Jones, Sato, Wada and Wiens wrote one down explicitly, of degree $25$ in $26$ variables. Such a polynomial needs more than one variable (Exercise~\ref{ex:primepoly}).

The heart of the theorem is that exponentiation, a function of exponential growth, can be defined by a polynomial. Davis, Putnam and Robinson had proved in 1961 that every recognisable set is \emph{exponential} Diophantine, defined by an equation that may also use $u^v$ with unknown $u,v$; so it was enough to show that $\{(m,k,n):m=k^n\}$ is Diophantine. The tool is Pell's equation (Tome~1, Chapter~8). Fix an integer $a\ge2$ and let $d=a^2-1$, which is not a square since $(a-1)^2<d<a^2$. Define $x_n=x_n(a)$ and $y_n=y_n(a)$ by
\[
x_n+y_n\sqrt d=(a+\sqrt d)^n\qquad(n\ge0),
\]
so that $x_0=1$, $y_0=0$, $x_1=a$, $y_1=1$, and $x_{n+1}=ax_n+dy_n$, $y_{n+1}=x_n+ay_n$.

\begin{lemma}\label{lem:pellsol}
The solutions $(x,y)\in\mathbb N^2$ of $x^2-dy^2=1$ are exactly the pairs $(x_n,y_n)$, $n\ge0$.
\end{lemma}

\begin{proof}
Since $(a+\sqrt d)(a-\sqrt d)=1$, taking conjugates gives $x_n-y_n\sqrt d=(a-\sqrt d)^n$, and multiplying, $x_n^2-dy_n^2=1$. Conversely let $u=x+y\sqrt d$ with $x,y\in\mathbb N$ and $x^2-dy^2=1$; then $u\ge1$. Let $\varepsilon=a+\sqrt d>1$ and choose $n\ge0$ with $\varepsilon^n\le u<\varepsilon^{n+1}$. Then $v=u\varepsilon^{-n}=u(a-\sqrt d)^n=X+Y\sqrt d$ with $X,Y\in\mathbb Z$, $X^2-dY^2=1$ and $1\le v<\varepsilon$, so $v^{-1}=X-Y\sqrt d\in(\varepsilon^{-1},1]$. Hence $Y\sqrt d=\frac{v-v^{-1}}2\in\bigl[0,\frac{\varepsilon-\varepsilon^{-1}}2\bigr)=[0,\sqrt d)$, since $\varepsilon^{-1}=a-\sqrt d$. So $Y=0$, $X=1$, $v=1$ and $u=\varepsilon^n$; as $\sqrt d$ is irrational, $(x,y)=(x_n,y_n)$.
\end{proof}

\begin{lemma}\label{lem:pellcong}
For all $n\ge0$: \textup{(a)} $y_n\equiv n\pmod{a-1}$; \textup{(b)} for every integer $k$, $x_n-(a-k)y_n\equiv k^n\pmod{2ak-k^2-1}$; \textup{(c)} $y_{n+1}\ge a^n$.
\end{lemma}

\begin{proof}
(a) By the binomial formula, $y_n=\sum_{j\text{ odd}}\binom nja^{n-j}d^{(j-1)/2}$. Modulo $a-1$, $a\equiv1$ and $d=(a-1)(a+1)\equiv0$, so only $j=1$ remains: $y_n\equiv n$. (b) Let $N=2ak-k^2-1$, so that $d=a^2-1\equiv(a-k)^2\pmod N$, and $z_n=x_n-(a-k)y_n$. Then $z_0=1$ and, by the recurrences,
\[
z_{n+1}=kx_n+\bigl(d-a(a-k)\bigr)y_n\equiv kx_n+\bigl((a-k)^2-a(a-k)\bigr)y_n=k\,z_n\pmod N.
\]
So $z_n\equiv k^n$. (c) $y_{n+1}=x_n+ay_n\ge ay_n$ and $y_1=1$.
\end{proof}

\begin{theorem}[Matiyasevich, 1970]\label{thm:matiyasevich}
The set $\{(a,y,n):a\ge2,\ y=y_n(a)\}$ is Diophantine.
\end{theorem}

This is the lemma that Robinson had been waiting for; we quote it, and refer to Davis's 1973 paper for a complete proof by elementary number theory, using the congruences above and the periodicity of $y_n$ modulo $y_m$. With it, exponentiation follows.

\begin{proposition}[Robinson]\label{prop:exponential}
The set $\{(m,k,n):m=k^n\}$ is Diophantine.
\end{proposition}

\begin{proof}
For $k=0$, $m=k^n$ means $(n=0\wedge m=1)$ or $(n\ge1\wedge m=0)$, a Diophantine condition, and we may combine cases with Proposition~\ref{prop:dioph}(b). Let $k\ge1$. We claim that $m=k^n$ if and only if there are $a,b,x,y,y'$ such that
\[
y'=y_{n+1}(k+1),\qquad a=y'+k+1+b,\qquad y=y_n(a),\qquad x^2-(a^2-1)y^2=1,
\]
\[
m\equiv x-(a-k)y\pmod{2ak-k^2-1},\qquad m<2ak-k^2-1.
\]
Each condition is Diophantine by Theorem~\ref{thm:matiyasevich} and Proposition~\ref{prop:dioph}, with a congruence $u\equiv v\pmod N$ written as $\exists z\ (u-v)^2=(Nz)^2$. Suppose they hold, and let $N=2ak-k^2-1$. As $a\ge2$, Lemma~\ref{lem:pellsol} forces $x=x_n(a)$, and Lemma~\ref{lem:pellcong}(b) gives $x-(a-k)y\equiv k^n$. Moreover $a\ge k+1$, so $N=k(2a-k)-1\ge k(a+1)-1\ge a$, and $a>y_{n+1}(k+1)\ge(k+1)^n>k^n$ by Lemma~\ref{lem:pellcong}(c) (with $k+1\ge2$ in place of $a$). So $0\le k^n<N$ and $0\le m<N$ are both the remainder of $x-(a-k)y$ modulo $N$, and $m=k^n$. Conversely, if $m=k^n$, all conditions hold with $b=0$.
\end{proof}

Together with the theorem of Davis, Putnam and Robinson, which we quote, Proposition~\ref{prop:exponential} gives Theorem~\ref{thm:mrdp}. The number of unknowns needed is surprisingly small: Matiyasevich and Robinson brought it to $13$ in 1975, and Matiyasevich to $9$, a proof written up by James Jones in 1982.

\section{Busy beavers and open problems}\label{sec:openbeavers}

The halting problem connects undecidability with the hardest open problems of mathematics. Goldbach's conjecture, that every even number $n\ge4$ is a sum of two primes, holds if and only if the machine that tests $n=4,6,8,\dots$ in turn and halts at the first counterexample runs forever; so does Fermat's Last Theorem for its own machine, and so would the Riemann hypothesis, by a more elaborate search. Each such machine has some number $n$ of states, and knowing $S(n)$, the busy beaver of Theorem~\ref{thm:beaver}, would settle the conjecture by running the machine for $S(n)+1$ steps. The busy beaver function concentrates the open problems of mathematics into its values, which is a vivid way to see why it grows so fast.

At six states the problems are no longer far away. In June 2024 the bbchallenge participant known as mxdys reported a six-state machine, soon named \emph{Antihydra}, whose behaviour reduces to a problem of the Collatz type: start with $a=8$ and a counter $b=0$, and repeatedly add $2$ to $b$ if $a$ is even and subtract $1$ if $a$ is odd, then replace $a$ by $a+\lfloor a/2\rfloor$; the machine halts if $b$ ever reaches $-1$, that is, if the odd values ever outnumber twice the even ones. Heuristically the parities behave like coin tosses and $b$ drifts upwards, so the machine should never halt, but nobody can prove it. In June 2025 mxdys found a six-state machine that halts after more than $2\uparrow\uparrow\bigl(2\uparrow\uparrow(2\uparrow\uparrow9)\bigr)$ steps, where $2\uparrow\uparrow n$ is a tower of $n$ twos: a lower bound for $S(6)$.

Further away, the questions cannot be settled at all from the usual axioms. A machine that searches for a contradiction in ZFC set theory runs forever if ZFC is consistent, but, by Gödel's second incompleteness theorem (Chapter~7), ZFC cannot prove it. Adam Yedidia and Scott Aaronson described in 2016 a machine with $7910$ states whose running forever cannot be proved in ZFC, relying on a theorem of Harvey Friedman; within weeks Stefan O'Rear built one with fewer than $750$ states that searches for a contradiction directly, and by 2026 enthusiasts had brought the number down to a few hundred. So the values of $S(n)$ are not even provable beyond some moderate $n$: the boundary between the knowable and the unknowable is somewhere between $6$ and a few hundred states.

\begin{beauty}
None of these problems was invented to be undecidable. Thue's words, Post's dominoes, Wang's tiles, Conway's fractions and Hilbert's equations came from algebra, logic, geometry and number theory, and each turned out to be a programming language in disguise: rewriting is computing, matching dominoes is checking a computation history, a tiled floor is a space-time diagram, and a polynomial equation with a solution is a computation that halts. Once a problem can simulate a universal machine, it inherits the halting problem, whatever it looks like. And the dividing line is low: one more letter, one more tile, one more state, and a problem that was a puzzle becomes as hard as all of mathematics.
\end{beauty}

\section{Exercises}

\begin{exercise}[$\star$]\label{ex:pcpunary}
Show that Post's correspondence problem is decidable for instances over a one-letter alphabet.
\end{exercise}

\begin{exercise}[$\star$]\label{ex:pcphand}
Find a match of the instance $(a,baa)$, $(ab,aa)$, $(bba,bb)$, and show that no shorter match exists.
\end{exercise}

\begin{exercise}[$\star$]\label{ex:commutative}
Show that the word problem of the semigroup presented by generators $a,b$ and the relation $ab=ba$ is decidable, and identify this semigroup.
\end{exercise}

\begin{exercise}[$\star\star$]\label{ex:singlyperiodic}
Show that if a tile set has a tiling $\tau$ with $\tau(x+p,y)=\tau(x,y)$ for all $x,y$ (one period), then it has a periodic tiling in the sense of Proposition~\ref{prop:wangperiodic}.
\end{exercise}

\begin{exercise}[$\star\star$]\label{ex:fractranmult}
Explain why the FRACTRAN program $\frac32$, started on $2^a3^b$, halts on $3^{a+b}$, and why $\frac{455}{33},\frac{11}{13},\frac1{11},\frac37,\frac{11}2,\frac13$, started on $2^a3^b$, halts on $5^{ab}$.
\end{exercise}

\begin{exercise}[$\star$]\label{ex:diophsets}
Show that the following sets are Diophantine: $\{(x,y,z):x\equiv y\pmod z\}$; the numbers that are not powers of $2$; the numbers that are not perfect squares.
\end{exercise}

\begin{exercise}[$\star\star$]\label{ex:integers}
(a) Show that $P(y_1,\dots,y_m)=0$ has a solution in $\mathbb N^m$ if and only if $P(s_1^2+t_1^2+u_1^2+v_1^2,\dots)=0$, with each $y_i$ replaced by a sum of four squares of new unknowns, has a solution in integers. (b) Show that $P=0$ has a solution in $\mathbb Z^m$ if and only if $\prod_{\varepsilon\in\{\pm1\}^m}P(\varepsilon_1y_1,\dots,\varepsilon_my_m)=0$ has a solution in $\mathbb N^m$. Conclude that Hilbert's tenth problem is unsolvable over $\mathbb Z$ as well as over $\mathbb N$.
\end{exercise}

\begin{exercise}[$\star\star$]\label{ex:pelladd}
Show that $x_{m+n}=x_mx_n+dy_my_n$ and $y_{m+n}=x_my_n+x_ny_m$, and deduce that $y_n$ divides $y_{nk}$ for all $n,k$.
\end{exercise}

\begin{exercise}[$\star$]\label{ex:primepoly}
Show that no nonconstant polynomial $P\in\mathbb Z[x]$ takes only prime values on $\mathbb N$.
\end{exercise}

\begin{exercise}[$\star\star\star$]\label{ex:twocounters}
Show that every counter machine can be simulated by one with two counters, started with the first counter equal to $2^{c_1}3^{c_2}\cdots p_k^{c_k}$ and the second equal to $0$. Deduce that the halting problem for counter machines with two counters is undecidable.
\end{exercise}

\section{Hints}

\begin{hint}{ex:pcpunary}
A domino $(a^p,a^q)$ changes the difference of lengths by $p-q$.
\end{hint}
\begin{hint}{ex:pcphand}
Which dominoes can start a match? Then follow the letter that the shorter row needs next.
\end{hint}
\begin{hint}{ex:commutative}
Count the letters, and sort a word.
\end{hint}
\begin{hint}{ex:singlyperiodic}
A tiling with period $(p,0)$ is determined by the $p$ tiles of each row in a strip of width $p$. Two rows of the strip are the same.
\end{hint}
\begin{hint}{ex:fractranmult}
Read each prime as a counter and the primes $11,13$ as instruction numbers, and translate the program into a counter machine.
\end{hint}
\begin{hint}{ex:diophsets}
$z$ divides $x-y$; a number is not a power of $2$ when it has an odd divisor greater than $1$; a non-square lies strictly between two consecutive squares.
\end{hint}
\begin{hint}{ex:integers}
Lagrange's four-square theorem (Tome~1, Chapter~7) for (a); for (b), a solution in $\mathbb Z$ becomes one in $\mathbb N$ after changing some signs.
\end{hint}
\begin{hint}{ex:pelladd}
Multiply $(x_m+y_m\sqrt d)(x_n+y_n\sqrt d)$, and expand $(x_n+y_n\sqrt d)^k$.
\end{hint}
\begin{hint}{ex:primepoly}
If $P(x_0)=p$ is prime, look at $P(x_0+jp)$ modulo $p$.
\end{hint}
\begin{hint}{ex:twocounters}
Multiplying by $p_i$, and testing divisibility by $p_i$, can be done by moving the first counter into the second, with the finite control counting up to $p_i$.
\end{hint}

\section{Solutions}

\begin{solution}{ex:pcpunary}
\idea Only lengths matter; a match balances the excesses of top over bottom.
\writeup Let the dominoes be $(a^{p_i},a^{q_i})$ and $d_i=p_i-q_i$. A sequence of indices is a match if and only if the sum of the $d_i$ over it is $0$. If some $d_i=0$, the single domino $i$ is a match. If $d_i>0$ and $d_j<0$ for some $i,j$, then $|d_j|$ copies of $i$ and $d_i$ copies of $j$ form a match. If all $d_i>0$, or all $d_i<0$, every nonempty sum is nonzero and there is no match. Checking these conditions is an algorithm.
\end{solution}

\begin{solution}{ex:pcphand}
\idea A match must begin with a domino whose two words agree on their common length; from there the choices are forced or quickly die.
\writeup Number the dominoes $1,2,3$. Domino $1$ starts with $a$ on top and $b$ at the bottom, and domino $2$ with $ab$ and $aa$; only domino $3$ can start, giving $bba$ over $bb$. The bottom must continue with $a$, so the next domino is $2$ (the only bottom starting with $a$): $bbaab$ over $bbaa$. The bottom must continue with $b$, which allows domino $1$ ($bbaaba$ over $bbaabaa$) or domino $3$ ($bbaabbba$ over $bbaabb$). After the second, the bottom must continue with $ba$, and domino $1$ gives
\[
bba\cdot ab\cdot bba\cdot a=bbaabbbaa=bb\cdot aa\cdot bb\cdot baa,
\]
the match $3,2,3,1$. No shorter match exists: a match starts with $3$, $3$ alone and $3,2$ are not matches, and the only sequences of length $3$ that can start a match are $3,2,1$ and $3,2,3$, whose rows ($bbaaba$ against $bbaabaa$, $bbaabbba$ against $bbaabb$) differ in length. (The notebook confirms by exhaustive search that $3,2,3,1$ is the only match of length at most $4$.)
\end{solution}

\begin{solution}{ex:commutative}
\idea The relation preserves the number of each letter, and allows us to sort any word.
\writeup Applying $ab\to ba$ or $ba\to ab$ does not change the numbers $|w|_a$ and $|w|_b$ of $a$'s and $b$'s. Conversely, every word can be turned into $a^{|w|_a}b^{|w|_b}$ by moving the $a$'s to the left one swap at a time. So $x=y$ in the semigroup if and only if $|x|_a=|y|_a$ and $|x|_b=|y|_b$, which is decidable. The map $w\mapsto(|w|_a,|w|_b)$ identifies the semigroup with $\mathbb N^2\setminus\{(0,0)\}$ under addition.
\end{solution}

\begin{solution}{ex:singlyperiodic}
\idea Look at the rows of a vertical strip of width $p$: there are finitely many possible rows, so one repeats, and the band between two equal rows can be stacked forever.
\writeup For each $y$ let $\rho(y)=(\tau(0,y),\dots,\tau(p-1,y))\in T^p$. There are finitely many such tuples, so $\rho(y_1)=\rho(y_2)$ for some $y_1<y_2$; let $q=y_2-y_1$. Define $\tau'(x,y)=\tau(x\bmod p,\ y_1+((y-y_1)\bmod q))$. Horizontal neighbours match: inside the strip because $\tau$ is a tiling, and across $x\equiv-1,0\pmod p$ because $\tau(p,y)=\tau(0,y)$. Vertical neighbours match inside the band $y_1\le y<y_2$ because $\tau$ is a tiling, and between the rows $y_2-1$ and $y_2\equiv y_1$ because the top colours of row $y_2-1$ match the bottom colours of $\rho(y_2)=\rho(y_1)$. So $\tau'$ is a tiling invariant under $(p,0)$ and $(0,q)$, hence under $(pq,0)$ and $(0,pq)$.
\end{solution}

\begin{solution}{ex:fractranmult}
\idea A FRACTRAN program is a counter machine whose counters are exponents; the primes in denominators that are not counters act as instruction numbers.
\writeup For $\frac32$: on $2^a3^b$ it applies while $a>0$, replacing a factor $2$ by $3$, and halts on $3^{a+b}$. For the second program, the exponents of $2,3,5,7$ are counters $a,b,r,t$, and $11,13$ mark two instructions. With no $11$ or $13$ present, the first three fractions do not apply. If $t>0$, $\frac37$ moves $t$ back into $b$. If $t=0$ and $a>0$, $\frac{11}2$ decrements $a$ and enters the instruction marked $11$. There, $\frac{455}{33}=\frac{5\cdot7\cdot13}{3\cdot11}$ applies while $b>0$: it decrements $b$, increments $r$ and $t$, and goes to $13$, from which $\frac{11}{13}$ returns to $11$. When $b=0$, $\frac1{11}$ leaves the instruction. So each unit of $a$ adds $b$ to $r$, copying $b$ into $t$ and back. When $a=0$ (and $t=0$), $\frac13$ erases $b$, and the program halts on $5^{ab}$.
\end{solution}

\begin{solution}{ex:diophsets}
\idea Write each condition with existential unknowns, then combine with Proposition~\ref{prop:dioph}(b).
\writeup $x\equiv y\pmod z$ if and only if $\exists u\ (x-y)^2-(zu)^2=0$. A number $x$ is not a power of $2$ if and only if $x=0$ or $x$ has an odd divisor $\ge3$: $x\cdot\bigl(x-u(2v+3)\bigr)=0$ for some $u,v$. A number $x$ is not a square if and only if $y^2<x<(y+1)^2$ for some $y$, that is, $\exists y,z,w$ with $(x-y^2-1-z)^2+\bigl((y+1)^2-x-1-w\bigr)^2=0$.
\end{solution}

\begin{solution}{ex:integers}
\idea Naturals are sums of four squares, and integers are naturals up to sign.
\writeup (a) If $\bar y\in\mathbb N^m$ is a solution, write each $y_i$ as a sum of four squares (Lagrange's theorem, Tome~1, Theorem~7.4) to get an integer solution of the new equation. Conversely, an integer solution of the new equation gives the solution $y_i=s_i^2+t_i^2+u_i^2+v_i^2\in\mathbb N$. (b) If $\bar y\in\mathbb Z^m$ solves $P=0$, then $(|y_1|,\dots,|y_m|)\in\mathbb N^m$ makes the factor with $\varepsilon_i=\operatorname{sign}y_i$ vanish. Conversely, if the product vanishes at $\bar z\in\mathbb N^m$, some factor does, and $(\varepsilon_1z_1,\dots,\varepsilon_mz_m)$ is an integer solution. So an algorithm for either version would give one for the other, and Corollary~\ref{cor:hilbert10} rules out both.
\end{solution}

\begin{solution}{ex:pelladd}
\idea Powers of $a+\sqrt d$ multiply; compare the coefficients of $\sqrt d$, which is irrational.
\writeup $(x_m+y_m\sqrt d)(x_n+y_n\sqrt d)=(x_mx_n+dy_my_n)+(x_my_n+x_ny_m)\sqrt d$, and this is $(a+\sqrt d)^{m+n}=x_{m+n}+y_{m+n}\sqrt d$; since $\sqrt d$ is irrational, the coefficients agree. Next, $x_{nk}+y_{nk}\sqrt d=(x_n+y_n\sqrt d)^k=\sum_j\binom kjx_n^{k-j}y_n^j(\sqrt d)^j$, and the coefficient of $\sqrt d$ collects the terms with $j$ odd, each divisible by $y_n$. So $y_n\mid y_{nk}$.
\end{solution}

\begin{solution}{ex:primepoly}
\idea A polynomial is periodic modulo $p$, so once it hits the prime $p$ it is divisible by $p$ infinitely often.
\writeup Suppose $P$ is nonconstant and all $P(x)$, $x\in\mathbb N$, are prime; let $p=P(0)$. For every $j$, $P(jp)\equiv P(0)\equiv0\pmod p$, because $jp\equiv0$ and $P$ has integer coefficients. So $P(jp)$ is a prime divisible by $p$, hence $P(jp)=p$ for all $j\ge0$. Then $P-p$ has infinitely many roots and is zero, and $P$ is constant, a contradiction.
\end{solution}

\begin{solution}{ex:twocounters}
\idea Store all counters as exponents of primes in one number, as FRACTRAN does, and use the second counter as scratch space for multiplying and dividing by a prime.
\writeup Let $A=\prod p_i^{c_i}$ and $Z=0$. To simulate $\mathrm{INC}(i,j)$, multiply $A$ by $p_i$: repeat ``decrement $A$, then increment $Z$ $p_i$ times'' until $A=0$; then repeat ``decrement $Z$, increment $A$'' until $Z=0$. To simulate $\mathrm{DEC}(i,j,l)$, divide $A$ by $p_i$: repeat ``decrement $A$ up to $p_i$ times, and increment $Z$ once if all $p_i$ succeed'', remembering in the instruction number the number $r<p_i$ of successful decrements in the last round, so that at the end $Z=\lfloor A/p_i\rfloor$ and $r=A\bmod p_i$. If $r=0$, $c_i>0$: move $Z$ back into $A$, which now codes $c_i-1$, and go to $j$. If $r>0$, $c_i=0$: rebuild $A=p_iZ+r$ by moving $Z$ into $A$ with $p_i$ increments per unit and $r$ more, and go to $l$. $\mathrm{HALT}$ stays $\mathrm{HALT}$. The two-counter machine halts if and only if the original one does. Applied to the three-counter machines of Theorem~\ref{thm:minsky}, started with counters $0,\mathrm{code}(w),0$, this gives two-counter machines started with $A=3^{\mathrm{code}(w)}$ and $Z=0$ that halt if and only if $M$ halts on $w$; so the halting problem for two counters is undecidable.
\end{solution}

\section*{Sources and further reading}

\begin{itemize}
\item \emph{A paper to read.} M.~Davis, \emph{Hilbert's tenth problem is unsolvable}, American Mathematical Monthly 80 (1973), 233--269. A complete, self-contained proof of the MRDP theorem, written for a wide audience. Read first the definitions and closure properties of Diophantine sets (our Proposition~\ref{prop:dioph}); then the section on the Pell equation, which proves Lemmas~\ref{lem:pellsol} and~\ref{lem:pellcong} and Theorem~\ref{thm:matiyasevich} through a chain of short lemmas; then the bounded universal quantifier, the step where the method of Davis, Putnam and Robinson enters; and last the coding of computations. The paper shows how far elementary number theory can go when every step is aimed at one goal.
\item \emph{The stories.} D.~Hilbert, \emph{Mathematical problems}, translated by M.~W.~Newson, Bulletin of the American Mathematical Society 8 (1902), 437--479, for the tenth problem; J.~Gray, \emph{The Hilbert Challenge}, Oxford University Press, 2000, for the Paris lecture; C.~Reid, \emph{The autobiography of Julia Robinson}, in D.~J.~Albers, G.~L.~Alexanderson, C.~Reid (eds.), \emph{More Mathematical People}, Harcourt Brace Jovanovich, 1990, for the birthday wish, the telephone call and the letter to Matiyasevich; Y.~Matiyasevich, \emph{My collaboration with Julia Robinson}, The Mathematical Intelligencer 14 (1992), 38--45.
\item A.~Thue, \emph{Probleme über Veränderungen von Zeichenreihen nach gegebenen Regeln}, Skrifter utgit av Videnskapsselskapet i Kristiania, I. Mat.-Naturv. Klasse, 1914, no.~10; E.~L.~Post, \emph{Recursive unsolvability of a problem of Thue}, Journal of Symbolic Logic 12 (1947), 1--11; A.~A.~Markov, \emph{On the impossibility of certain algorithms in the theory of associative systems}, Doklady Akademii Nauk SSSR 55 (1947), 583--586; G.~S.~Tseitin, \emph{An associative calculus with an insoluble problem of equivalence}, Trudy Matematicheskogo Instituta im.\ V.~A.~Steklova 52 (1958), 172--189.
\item E.~L.~Post, \emph{A variant of a recursively unsolvable problem}, Bulletin of the American Mathematical Society 52 (1946), 264--268; the proof of Theorem~\ref{thm:pcp} through the modified problem follows M.~Sipser's textbook (Chapter~\ref{ch:automata}), with rewriting systems in place of machines as in J.~E.~Hopcroft and J.~D.~Ullman, \emph{Introduction to Automata Theory, Languages, and Computation}, Addison-Wesley, 1979.
\item H.~Wang, \emph{Proving theorems by pattern recognition II}, Bell System Technical Journal 40 (1961), 1--41; R.~Berger, \emph{The undecidability of the domino problem}, Memoirs of the American Mathematical Society 66 (1966); R.~M.~Robinson, \emph{Undecidability and nonperiodicity for tilings of the plane}, Inventiones Mathematicae 12 (1971), 177--209; E.~Jeandel and M.~Rao, \emph{An aperiodic set of 11 Wang tiles}, Advances in Combinatorics (2021), arXiv:1506.06492; D.~Smith, J.~S.~Myers, C.~S.~Kaplan and C.~Goodman-Strauss, \emph{An aperiodic monotile}, Combinatorial Theory 4 (2024).
\item M.~L.~Minsky, \emph{Recursive unsolvability of Post's problem of ``Tag'' and other topics in theory of Turing machines}, Annals of Mathematics 74 (1961), 437--455, and \emph{Computation: Finite and Infinite Machines}, Prentice-Hall, 1967; J.~H.~Conway, \emph{Unpredictable iterations}, Proceedings of the 1972 Number Theory Conference, University of Colorado, Boulder, 1972, 49--52; J.~H.~Conway, \emph{FRACTRAN: a simple universal programming language for arithmetic}, in T.~M.~Cover and B.~Gopinath (eds.), \emph{Open Problems in Communication and Computation}, Springer, 1987, 4--26; R.~K.~Guy, \emph{Conway's prime producing machine}, Mathematics Magazine 56 (1983), 26--33; S.~A.~Kurtz and J.~Simon, \emph{The undecidability of the generalized Collatz problem}, in \emph{Theory and Applications of Models of Computation}, Lecture Notes in Computer Science 4484, Springer, 2007, 542--553; J.~C.~Lagarias (ed.), \emph{The Ultimate Challenge: The $3x+1$ Problem}, American Mathematical Society, 2010.
\item M.~Davis, H.~Putnam and J.~Robinson, \emph{The decision problem for exponential Diophantine equations}, Annals of Mathematics 74 (1961), 425--436; Y.~V.~Matiyasevich, \emph{Enumerable sets are Diophantine}, Doklady Akademii Nauk SSSR 191 (1970), 279--282; J.~P.~Jones, D.~Sato, H.~Wada and D.~Wiens, \emph{Diophantine representation of the set of prime numbers}, American Mathematical Monthly 83 (1976), 449--464; J.~P.~Jones, \emph{Universal Diophantine equation}, Journal of Symbolic Logic 47 (1982), 549--571; Y.~V.~Matiyasevich, \emph{Hilbert's Tenth Problem}, MIT Press, 1993.
\item The bbchallenge Collaboration, pages on Antihydra and on logical independence at \texttt{bbchallenge.org} and at the bbchallenge wiki; S.~Aaronson, \emph{BusyBeaver(6) is really quite large}, blog post, June 2025; A.~Yedidia and S.~Aaronson, \emph{A relatively small Turing machine whose behavior is independent of set theory}, Complex Systems 25 (2016), 297--327.
\end{itemize}
